\documentclass[11pt,letterpaper]{article}

\usepackage[T1]{fontenc}
\usepackage[utf8]{inputenc}
\usepackage{times}
\usepackage[top=0.9in, bottom=0.9in, left=1in, right=1in]{geometry}
\usepackage{microtype}
\usepackage{booktabs}
\usepackage{amsmath}
\usepackage{amssymb}
\usepackage{url}
\usepackage[hidelinks]{hyperref}
\usepackage{parskip}
\usepackage{float}

\makeatletter
\renewcommand{\section}{\@startsection{section}{1}{\z@}%
  {5pt \@plus 1pt \@minus 1pt}{1.5pt}%
  {\normalfont\normalsize\bfseries}}
\renewcommand{\subsection}{\@startsection{subsection}{2}{\z@}%
  {3pt \@plus 1pt}{1pt}%
  {\normalfont\normalsize\itshape}}
\makeatother
\setlength{\parskip}{2.5pt}
\setlength{\parindent}{0pt}

\title{\vspace{-1.5em}%
  \textbf{CS505 Project Status Report: Citation-Grounded QA with Local Models}\\[3pt]
  {\normalsize\normalfont
   Boston University \enspace|\enspace
   Team Progress Report \enspace|\enspace
   April 8, 2026}\\[1pt]
  {\normalsize\normalfont Jaswanth Pinnepu | Zukhriddin Fakhriddinov | Atharv Gulati | Aayush Kumar}%
}
\date{\vspace{-1.5em}}

\begin{document}
\maketitle
\thispagestyle{empty}

\section{Project Goal}

The project studies hallucination reduction in scientific question answering using a fully local retrieval-augmented generation pipeline. The current stack is built around QASPER, shared SCC storage, and locally-run small instruction-tuned models. The main experimental question is how much answer quality and hallucination behavior depend on retrieval quality versus generation behavior, especially under citation forcing.

\section{SCC Layout and Reproducibility}

The implementation uses a strict split between shared deterministic artifacts and per-user experimental outputs.

\textbf{Shared course path}

\texttt{/projectnb/cs505am/projects/zukhriddin\_nlp\_research\_project/qasper}

\textbf{Student working repo}

\texttt{/projectnb/cs505am/students/zfmsai/zukhriddin\_nlp\_research\_project}

\textbf{Backward-compatible symlink}

\texttt{/projectnb/cs505am/students/zfmsai/nlp\_research\_project}

The code automatically resolves either the renamed project path or the older path. Shared data and processed artifacts live in \texttt{projects/}; model caches, runs, and temporary indexes live in \texttt{students/\$USER/}.

\section{What Was Implemented}

\subsection{Task 2: Data and Chunking}

Implemented components:

\begin{itemize}
\item QASPER loader from official JSON/Hugging Face format
\item Internal normalized schema for papers, questions, answers, and chunks
\item Fixed chunking
\item Sliding-window chunking
\item Section-aware chunking
\item Deterministic processed split builder for team-shared reuse
\end{itemize}

Canonical shared processed artifacts were built for:

\begin{itemize}
\item \texttt{train}
\item \texttt{validation}
\item \texttt{test}
\end{itemize}

These outputs are the interface boundary used by both retrieval and generation.

\subsection{Task 3: Retrieval}

Implemented retrieval stack:

\begin{itemize}
\item BM25
\item Dense retrieval using \texttt{BAAI/bge-small-en-v1.5}
\item Hybrid reciprocal-rank fusion
\item Cross-encoder reranking using \texttt{cross-encoder/ms-marco-MiniLM-L-6-v2}
\item Evaluation via Recall@5, MRR@10, and evidence hit rate
\end{itemize}

\subsection{Task 4: Generation and Evaluation}

Implemented generation stack:

\begin{itemize}
\item Baseline prompt style
\item Citation-forcing prompt style
\item Local generation with \texttt{microsoft/Phi-3.5-mini-instruct} via \texttt{transformers}
\item Answer normalization for canonical QASPER answer forms (\texttt{yes}, \texttt{no}, \texttt{unanswerable})
\item Context compression to question-relevant sentences
\item SCC-safe batch scripts and memory settings
\item Evaluation via token F1, citation precision, citation rate, hallucination proxy, and supported-answer rate
\end{itemize}

\subsection{Task 1: Infra and SCC}

Implemented infrastructure:

\begin{itemize}
\item Reproducible path/config system
\item SCC batch scripts for retrieval and generation
\item Student-local Hugging Face cache handling
\item Documentation for SCC environment setup
\item GPU-safe defaults that run on \texttt{academic-gpu}
\end{itemize}

\section{How the Pipeline Works}

\begin{enumerate}
\item Load QASPER and normalize into internal objects.
\item Build processed split artifacts shared across the team.
\item Chunk each paper with one of the chunking strategies.
\item Retrieve from chunks belonging only to the paper for the given question.
\item Generate an answer from the top retrieved chunks.
\item Score the generated answer against gold answers and gold evidence.
\end{enumerate}

The current best retriever for generation is \texttt{hybrid}; \texttt{hybrid\_rerank} was evaluated later as a stronger comparator.

\section{Main Results}

\subsection{Retrieval Results on Validation (Section Chunks)}

\begin{table}[H]
\centering
\small
\begin{tabular}{@{}lccc@{}}
\toprule
\textbf{Method} & \textbf{Recall@5} & \textbf{MRR@10} & \textbf{Evidence Hit} \\
\midrule
BM25 & 0.5769 & 0.4110 & 0.7152 \\
Dense & 0.5886 & 0.4820 & 0.7238 \\
Hybrid & 0.6313 & 0.5097 & 0.7778 \\
Hybrid + rerank & 0.6410 & 0.5027 & 0.7702 \\
\bottomrule
\end{tabular}
\caption{Validation retrieval results on section-aware chunks.}
\end{table}

\textbf{Takeaway:} Hybrid retrieval is the strongest overall default for generation; hybrid + rerank improves recall slightly but did not clearly dominate all retrieval metrics.

\subsection{Generation Results with Hybrid Retrieval}

\begin{table}[H]
\centering
\small
\begin{tabular}{@{}lccccc@{}}
\toprule
\textbf{Prompt} & \textbf{Questions} & \textbf{Token F1} & \textbf{Citation Prec.} & \textbf{Halluc.} & \textbf{Supported} \\
\midrule
Baseline & 1005 & 0.3187 & 0.0000 & 0.4657 & 0.7612 \\
Citation forcing & 1005 & 0.3167 & 0.3423 & 0.5572 & 0.5741 \\
\bottomrule
\end{tabular}
\caption{Full-validation generation results using hybrid retrieval. Citation forcing had citation rate 0.9254.}
\end{table}

\textbf{Takeaway:} Citation forcing now works much better than the initial versions and produces citations reliably, but baseline remains better on answer quality, hallucination proxy, and supported-answer rate.

\subsection{Generation Results with Hybrid + Rerank}

\begin{table}[H]
\centering
\small
\begin{tabular}{@{}lccccc@{}}
\toprule
\textbf{Prompt} & \textbf{Questions} & \textbf{Token F1} & \textbf{Citation Prec.} & \textbf{Halluc.} & \textbf{Supported} \\
\midrule
Baseline & 100 & 0.3697 & 0.0000 & 0.4300 & 0.8000 \\
Citation forcing & 100 & 0.3569 & 0.3528 & 0.5400 & 0.6600 \\
\bottomrule
\end{tabular}
\caption{100-question validation sample using hybrid + rerank. Citation forcing had citation rate 0.9100.}
\end{table}

\textbf{Takeaway:} The strongest generation setting observed so far is the \texttt{hybrid\_rerank} baseline. Even under stronger retrieval, citation forcing still trades off answer quality for citation behavior.

\section{What Changed During Debugging}

Several implementation issues were identified and fixed during the experiments:

\begin{itemize}
\item SCC scheduler instability prevented job submission for a period; jobs were moved to proper batch usage once SCC recovered.
\item The correct GPU queue for this course environment was \texttt{academic-gpu}.
\item Generation initially hit CUDA OOM; the fix was to cap input/context lengths and use conservative CUDA settings.
\item Token F1 was artificially low because the model often produced verbose forms like ``The context does not specify...'' instead of exact canonical answers. Answer normalization fixed this.
\item SCC environment-variable submission did not reliably propagate \texttt{QUESTION\_LIMIT}; the batch script was changed to accept it as a positional argument.
\item After the SCC project rename, shared path resolution failed; config detection was extended to support both old and new project names.
\end{itemize}

\section{Current Status}

Completed:

\begin{itemize}
\item Task 2: complete
\item Task 3: complete
\item Task 4 core Phi-based evaluation pipeline: complete
\item Task 1 setup, SCC usage, and reproducibility documentation: complete for the current stack
\end{itemize}

Reached goals:

\begin{itemize}
\item deterministic shared QASPER preprocessing for the whole team
\item full retrieval implementation and benchmarking
\item full generation implementation and benchmarking with Phi-3.5
\item SCC-compatible batch execution for both retrieval and generation
\item clear best current setting before moving on to new models
\end{itemize}

\textbf{Best current setting before Qwen:}

\begin{itemize}
\item retrieval: \texttt{hybrid\_rerank} is the strongest generation-side comparator on the current 100-question sample
\item generation: baseline prompting is still better than citation forcing for raw answer quality
\end{itemize}

\section{How to Inspect Results}

From SCC:

\begin{verbatim}
cd /projectnb/cs505am/students/$USER/zukhriddin_nlp_research_project
ls -1t runs/retrieval/validation/section/*.json | head
ls -1t runs/generation/validation/section/*.json | head
\end{verbatim}

To print metrics from a saved run:

\begin{verbatim}
python - <<'PY'
import json
from pathlib import Path
run_dir = Path("runs/generation/validation/section")
path = run_dir / "hybrid.baseline.20260409T011901Z.json"
obj = json.loads(path.read_text())
print(obj["metrics"])
PY
\end{verbatim}

\section{How to Reproduce the Pipeline}

\subsection{Environment}

\begin{verbatim}
chmod +x setup_env.sh
./setup_env.sh
\end{verbatim}

\subsection{Check resolved paths}

\begin{verbatim}
python3 scripts/show_project_paths.py
\end{verbatim}

\subsection{Retrieval evaluation}

\begin{verbatim}
python scripts/evaluate_qasper_retrieval.py \
  --split validation \
  --strategy section \
  --method hybrid
\end{verbatim}

\subsection{Generation evaluation}

\begin{verbatim}
python scripts/evaluate_qasper_generation.py \
  --split validation \
  --strategy section \
  --retrieval-method hybrid \
  --prompt-style baseline \
  --generation-model microsoft/Phi-3.5-mini-instruct
\end{verbatim}

\subsection{SCC batch generation sample}

\begin{verbatim}
qsub -N qasper_gen_sample -o logs/qasper_gen_sample.log \
  scripts/qsub_qasper_generation.sh hybrid citation_forcing \
  validation section 5 20 100 2048 1200 128
\end{verbatim}

The final four positional arguments are:

\begin{itemize}
\item question limit
\item max input tokens
\item max context tokens
\item max new tokens
\end{itemize}

\section{Next Step}

The current Phi-based milestone is complete enough to move on. The next experiments should compare:

\begin{itemize}
\item \texttt{Qwen2.5-7B-Instruct}
\item \texttt{Qwen3-4B-Instruct-2507}
\end{itemize}

Those runs should reuse the same retrieval and evaluation pipeline so that model changes are isolated from data and retrieval changes.

\end{document}
